\scanpage{81}
\diagram{\begin{tikzcd}[ampersand replacement=\&]\unsure{\mathbb E|\mathrm E}=Y[t_1,\ldots t_{\unsure{r}}]\arrow[d,no head]\\Y\end{tikzcd}}
On peut supposer $Y$ intègre, \inserted{de dim $\unsure{n|\varphi}$,} $X$
plongé dans $\mathbb E\mathrel{\unsure{:|=}}Y[t_1,\ldots,t_r]$, $X$ intègre.
\deleted{Alors codim} Je dis que
\[
\operatorname{Codim}(X^{\unsure{\prime}},E)\geq n+r-\unsure{\delta|d|\nu}\quad\text{i.e. }\delta\geq\inserted{n+r-}\operatorname{codim}(X,E)
\]
En effet, si $X\ne X_0$, alors le premier membre est
égal à $\operatorname{codim}(X',E')\unsure{=}\dim E'-\dim X'$
% slanted
\margin{\unsure{tt|H} sur un ``bon'' Jacobson $Y'$}
\[
=((n\unsure{-}1)+r)\deleted{\unsure{+}}-\dim X'\geq(n-1+r)-(\delta-1)=n+r-\delta\quad\text{OK}
\]
Si $X=X_0$, alors $\delta$\deleted{$\unsure{\geq\sup(}=\dim X_0-$}
\[
\begin{aligned}
\operatorname{Codim}(X,E)&=\operatorname{Codim}(X_0,E_0)+\operatorname{Codim}(E_0,E)=\\
&=(r-\delta)+n\qquad\text{OK.}
\end{aligned}
\]
[NB. On \inserted{ne} peut avoir $\operatorname{codim}(X,E)>n+r-\delta$,
\deleted{sauf} \inserted{que} lorsque $\dim X_0>\dim X'+1$ \deleted{\unsure{pour}} \inserted{mais}
\deleted{\unsure{ou}} $X\ne X_0$.]
\margin{\uline{est-ce possible ?} \uuline{non}}
\deleted{Or, puisque $X_0\ne\varnothing$, \unsure{de} utilisant la formule des dimensions,}
\inserted{\unsure{Non.|N.B.} prouvons en fait}
\margin{\deleted{\unsure{$H^i$}} $\operatorname{cd}_{\ell}(X)\leq(n+r)-\operatorname{codim}(X,E)$}

On va procéder par récurrence sur $r$.
Pour $r=0$, c'est trivial \deleted{(tt est [\unsure{str}} (alors $X\subset Y$, donc $H^i(X,F)=0$ pour $i>0$ i.e.
$\operatorname{cd}(X)=0\leq\deleted{\unsure{r-n}}n\unsure{-|\geq}\operatorname{codim}(X,Y)=\dim X$. Si
$r=1$, c'est le th. $1_n$. Dans le cas
\deleted{général}, \inserted{$r\geq2$} on utilise la suite spectrale
pour $E_Y^r\longrightarrow E_Y^{r-1}$ \ldots\ldots
% archive sheet number: 80
